\title{Parameter-Efficient Orthogonal Finetuning via Factorization}
Drawing inspiration from the Cooley-Tukey fast Fourier transform, which enables efficient communication of information, we introduce an efficient orthogonal parameterization based on butterfly structures. To fully approximate all orthogonal matrices of size $d$, it suffices to compose products of butterfly matrices as $\bm{B}_{d-1,1}(d)\bm{B}^\top_{d-1,2}(d)\cdots\bm{B}_{1,1}(d)\bm{B}^\top_{1,2}(d)$, with $\bm{B}_{i,j}(d),\forall i,\forall j$ denoting butterfly matrices. \textbf{Spectral properties}. For instance, ensuring that each $\bm{B}_i$ is full-rank (which is required for orthogonality) necessitates at least $d$ edges to establish a bijective correspondence between nodes at the $i$-th level and nodes at the $(i=1)$-th level; this is the case in Figure~\ref{fig:info_trans}, where the minimum is 4 edges. \textbf{Butterfly structures}. The orthogonal transformation $\bm{R}(m,b)\in\mathbb{R}^{d\times d}$ is built as a sequence of products involving multiple orthogonal butterfly structures. Likewise, BOFT configurations with smaller $b$ and larger $m$ tend to offer improved parameter efficiency. In general, for a dense matrix $\bm{R}\in\mathbb{R}^{d\times d}$, the factor matrices $\bm{B}_m,\cdots,\bm{B}_1$ must collectively produce a system of directed edges over a $d\times (m+1)$ grid, where each edge links two nodes on successive layers (\emph{i.e.}, consecutive columns), so that any node at the first layer can reach every node at layer $(m+1)$. For efficient use of parameters, a block-diagonal construction is adopted, with $\bm{R}$ parameterized as $\text{diag}(\bm{R}_1,\bm{R}_2,\cdots,\bm{R}_r)$, where each $\bm{R}_i\in\mathbb{R}^{b\times b}$ is a small orthogonal matrix, and $br=d$. Overall, BOFT($m$,$b$) introduces $\frac{1}{2}(b-1)dm$ trainable parameters to adjust a linear transformation of shape $d\times n$. Using this framework, we identify the butterfly structure as a practical mechanism for achieving sparse orthogonal matrix decomposition. Furthermore, the sparsity characteristics associated with matrix factorization in BOFT may introduce additional forms of implicit inductive bias~\cite{gunasekar2017implicit,li2018algorithmic,liu2019neural}. Owing to its straightforward nature, this form of low-rank matrix reparameterization has become prevalent~\cite{dettmers2023qlora,zi2023delta,chavan2023one}. For brevity, we refer to BOFT that utilizes $\bm{R}(m,\frac{b}{2})$ as BOFT($m$,$b$) for $b\geq 2$. The underlying matrix factorization further allows for meaningful weight interpolation, as visualized in Figure~\ref{fig:boft_interpolation}. We then conduct comprehensive experiments to evaluate the adaptation of large vision transformers, extensive language models, and text-to-image diffusion models across a diverse set of downstream vision and language tasks. The core concept underlying model finetuning is to maintain a level of similarity between the pretrained and the finetuned model, as measured by certain criteria, thereby retaining knowledge acquired during pretraining. For instance, BOFT($9$,$2$) achieves higher expressivity than BOFT($1$,$16$) while requiring far fewer parameters. Here, each $\bm{B}^F_i(k)$ for $i=1,\dots,d/k$ corresponds to a butterfly factor as specified in Equation~\ref{eq:but_factor}. The butterfly architecture combined with permutations is theoretically capable of exactly representing a range of well-known fast linear transforms~\cite{dao2019learning,dao2019kaleidoscope} (\emph{e.g.}, the fast Fourier transform and Hadamard transform). The outcomes presented in Figure~\ref{fig:para_eff} are averages over 10 different random seeds. Specifically, we evaluate BOFT on multiple adaptation benchmarks spanning both computer vision and natural language processing. Our strategy for parameter-efficient information transmission draws upon the butterfly structures found in the Cooley-Tukey fast Fourier transform algorithm~\cite{cooley1965algorithm}, known for their efficiency in facilitating node-to-node communication~\cite{klasing1994broadcasting}. Yet, this improvement in parameter efficiency entails certain compromises: the orthogonal matrices adopt a predetermined sparsity pattern and transformations are executed independently within each block. Special cases include BOFT($1$,$b$), which collapses to the block-diagonal OFT~\cite{qiu2023controlling} having blocks of size $b$. The only requirement is that every factor $\tilde{\bm{B}}(d,k)$ in the composition of $\bm{B}(d)$ remains orthogonal. This strategy helps identify a more compact subset of hypotheses within the orthogonal group for use in downstream applications. OFT achieves this by learning an orthogonal transformation for the neurons within each layer, yielding both superior generalization and more stable optimization relative to LoRA. To contextualize this factorization process, we reinterpret the task of generating a dense orthogonal matrix as one of transmitting information through a network. While LoRA attains parameter-efficiency by adopting a low-rank approach, the original OFT instead employs a block-diagonal orthogonal matrix structure~\cite{qiu2023controlling}; reducing the size of the diagonal blocks increases parameter-efficiency within OFT. \textbf{Comparison to butterfly-based sparse training}. This expectation is grounded in the fact that the butterfly structure is uniquely capable of realizing many foundational linear transformations, a property that the block-diagonal approach in OFT does not possess. For the case where $d=2^N$, the $d$-dimensional \emph{butterfly matrix} $\bm{B}(d) \in \mathbb{R}^{d\times d}$ is constructed recursively as

\begin{equation}
\begin{aligned}
    \!\!\bm{B}(d)=\tilde{\bm{B}}(d,d)\!\cdot\!\begin{bmatrix}
    \bm{B}_1(\frac{d}{2})\!\!\!\!\! The information transmission viewpoint adopted here opens the possibility of leveraging insights from computer networking, a field extensively concerned with efficient network topology design. This observation naturally leads to an important question: \emph{Is it possible to construct a dense orthogonal matrix that remains parameter-efficient?}

In response to this question, we introduce a method that constructs a dense orthogonal matrix through the multiplication of several sparse orthogonal matrices. This can be illustrated using singular value decomposition: express $\bm{W}^0$ as $\bm{U}\bm{\Sigma}\bm{V}^\top$, where $\bm{U}$ and $\bm{V}$ are orthogonal matrices, and $\bm{\Sigma}$ is diagonal, containing the singular values. Such a block-diagonal configuration, while reducing parameter overhead, may induce less desirable inductive biases (\emph{e.g.}, diminished representational capacity and inability to approximate general linear transformations), and crucially, the optimal selection of a sparsity pattern remains an unresolved issue. Numerous approaches to PEFT exist~\cite{houlsby2019parameter,aghajanyan2020intrinsic,hulora2022,edalati2022krona,wang2022adamix,gheini2021cross,zaken2022bitfit,guo2020parameter,sung2021training,ansell2022composable,lester2021power,li2021prefix,vu2022spot,he2021towards,mao2021unipelt,karimi2021compacter,liu2022few,sung2022lst,chen2022parameter,jia2022visual,chen2022adaptformer,zhang2022neural,jie2023fact,lian2022scaling,luo2023towards}; among these, reparameterization-based techniques~\cite{aghajanyan2020intrinsic,hulora2022,edalati2022krona} are especially pertinent to our contribution. Similarly, for higher block sizes $k>2$, the structure of $\tilde{\bm{B}}^b(d,k)$ corresponds to a block-wise permutation of the pattern in $\tilde{\bm{B}}^b(d,2)$, so these can also be parameterized to maintain orthogonality. This graphical framework permits an exploration of OFT's parameter efficiency, and it readily facilitates the design of efficient matrix factorizations. The rationale for utilizing orthogonal transformations during the finetuning of weights lies in maintaining the angles between neuron pairs~\cite{liu2018learning,liu2021orthogonal,qiu2023controlling}, thus preserving much of the semantic knowledge acquired during pretraining. Our proposed multiplicative Dropout procedure randomly selects a portion $p_1$ of butterfly modules and a portion $p_2$ of their diagonal blocks, replacing the chosen blocks by identity matrices. This methodology is motivated by the realization that a trade-off exists between memory usage and computational effort: parameter count can be minimized at the expense of increased computation during training, as the product of sparse matrices must be repeatedly calculated each iteration. To ensure zero initialization, $\bm{R}$ is initialized as the identity. Such matrices have been leveraged for representing orthogonal matrices without requiring pivoting in Gaussian elimination~\cite{parker1995random}, enhancing computational efficiency, stabilizing recurrent neural network training~\cite{jing2017tunable}, and approximating kernels~\cite{munkhoeva2018quadrature}. Central to our method is the realization that leveraging a product of multiple sparse orthogonal matrices enables an efficient parameterization of a full dense orthogonal matrix. As another example, BOFT($6$,$4$) achieves a comparable approximation error to BOFT($2$,$16$), yet utilizes significantly fewer parameters. It is instructive to compare our sparsity-driven design with the low-rank methodology used in LoRA~\cite{hulora2022}; both low-rank and sparse matrices are prominent categories of structured matrices~\cite{candes2011robust} renowned for their parameter efficiency. Nonetheless, the practical deployment of OFT is challenged by the considerable number of parameters stemming from the large size of orthogonal matrices. This innovation allows us to extend orthogonal finetuning to diverse adaptation problems. Specifically, OFT trains an orthogonal matrix $\bm{R}\in\mathbb{R}^{d\times d}$ for a pretrained linear transformation $\bm{W}^0\in\mathbb{R}^{d\times n}$, altering the forward computation from $\bm{z}=(\bm{W}^0)^\top\bm{x}$ to $\bm{z}=(\bm{R}\bm{W}^0)^\top\bm{x}$, where $\bm{x}\in\mathbb{R}^{d}$ represents the input and $\bm{z}\in\mathbb{R}^{n}$ the output. \end{itemize}

\textbf{Parameter-efficient finetuning (PEFT)}. As a specific illustration (see Figure~\ref{fig:info_trans}), consider the decomposition $\bm{R}=\bm{B}_5\bm{B}_4\bm{B}_3\bm{B}_2\bm{B}_1$; the associated sparsity patterns and resultant graph structure are visualized. These $d$ edges are essential for orthogonality and are not included in the edge count, as these entries are not adjustable during training (for instance, in a $d\times d$ orthogonal matrix with $d$ nonzero entries, each can only be $\pm 1$). In addition, butterfly structures are effective in fast matrix-vector products~\cite{michielssen1996multilevel,o2010algorithm}, in approximating data-sparse matrices~\cite{li2015butterfly}, and in various aspects of network transmission~\cite{klasing1994broadcasting,li2003linear,rajasingh2016transmission}. In summary, the forward computation in BOFT is given by

\begin{equation}
    \begin{aligned}\nonumber
    \bm{z}=\big(\bm{R}(m,b)\cdot\bm{W}^0\big)^\top\bm{x},~~~~\text{s.t.}~\bigg{\{}\bm{R}(m,b)=\prod_{i=1}^m \tilde{\bm{B}}^b_{(d,i)}~~\&~~\underbrace{\big(\tilde{\bm{B}}^b_{(d,j)}\big)^\top\tilde{\bm{B}}^b_{(d,j)}=\tilde{\bm{B}}^b_{(d,j)}\big(\tilde{\bm{B}}^b_{(d,j)}\big)^\top\!=\bm{I}_d}_{\forall j\in [1, m]}\bigg{\}},
    \end{aligned}
\end{equation}

In this formulation, for notational simplicity, $\tilde{\bm{B}}^b(d,2^{m - i+1})$ is written as $\tilde{\bm{B}}^b_{(d,i)}$, and $\bm{I}_d$ stands for the $d$-dimensional identity matrix. Furthermore, the extent to which OFT generalizes to a broader range of adaptation scenarios (outside its original context in text-to-image diffusion~\cite{qiu2023controlling}) remains an open question. This work examines a rigorous adaptation strategy known as Orthogonal Finetuning (OFT) for tuning foundation models to downstream applications. Our objective is to build a dense orthogonal matrix by composing $m$ sparse orthogonal matrices, aiming to use the fewest effective trainable parameters possible. \textbf{BOFT Identity Initialization}. BOFT encompasses OFT as a particular instance, thus offering a broader orthogonal finetuning framework. Drawing on insights from hyperspherical energy and its relationship to generalization~\cite{liu2018learning,liu2021orthogonal}, \cite{qiu2023controlling} introduce orthogonal finetuning (OFT), which provides an alternative strategy for updating text-to-image diffusion models. To establish a coherent and accessible framework for examining the sparsity structures in orthogonal matrix decompositions, we reinterpret the task of constructing a dense orthogonal matrix in OFT as an analogy to information transmission. Exploiting this parameter-efficient orthogonalization, we introduce Orthogonal Butterfly~(BOFT), a new finetuning method that generalizes and extends OFT. Furthermore, it remains uncertain whether butterfly networks represent the most optimal structure for information transfer. The work of \cite{dao2022monarch} explores fine-tuning by first projecting pretrained weights onto a butterfly matrix variant, subsequently updating only the projected elements for target applications. This method constructs a dense matrix of dimension $d$ through the product of $\mathcal{O}(\log d)$ sparse matrices, where each individual sparse matrix contains only $\mathcal{O}(d)$ non-zero elements. The vertical axis measures the approximation error, while the horizontal axis indicates the number of effective trainable parameters. To mitigate this issue, we first analyze OFT through the lens of information transmission, leading us to identify several essential properties that underpin greater parameter efficiency. In essence, methods for the efficient adaptation of existing foundation models to task-specific settings are vital. Nevertheless, the number of trainable parameters in OFT typically exceeds that in LoRA. \end{theorem}

Theorem~\ref{thm:exp_boft} leads to a natural extension of BOFT—namely, generalizing the target orthogonal matrix as $\thickmuskip=2mu \medmuskip=2mu \bm{R}^G(m_1,b_1,m_2,b_2,l)=\bm{R}_{l,1}(m_1,b_1)\bm{R}_{l,2}^T(m_2,b_2)\cdots\bm{R}_{1,1}(m_1,b_1)\bm{R}_{1,2}^T(m_2,b_2)$, where $\bm{R}_{i,j}^T(m,b)$ represents the orthogonal matrix employed in BOFT. The \emph{butterfly component} is defined as $\thickmuskip=2mu \medmuskip=2mu \tilde{\bm{B}}(d,k)=\text{diag}(\bm{B}^F_1(k),\cdots,\bm{B}^F_{d/k}(k))$, forming a block-diagonal matrix of size $d\times d$ with each block of size $k$. This approach effectively represents orthogonal matrices using compositions of several $2\times 2$ orthogonal matrices. In contrast to the block-diagonal formulation employed in OFT, which attempts to balance expressiveness and structural regularity, BOFT leverages the butterfly structure to facilitate a more continuous spectrum between full orthogonal matrices (maximal expressiveness) and identity matrices (maximal regularity). The next section discusses how the butterfly architecture contributes to improved parameter efficiency. To mitigate this, OFT employs a block-diagonal formulation to decrease the parameter count. Maintaining the orthogonality of $\bm{R}$ during training is achieved through Cayley parameterization as described in \cite{qiu2023controlling,liu2021orthogonal}: $\bm{R}=(\bm{I}+\bm{Q})(\bm{I}-\bm{Q})^{-1}$, where $\bm{Q}$ is skew-symmetric ($\bm{Q}=-\bm{Q}^\top$). This perspective arises from viewing a dense $d\times d$ matrix as facilitating full connectivity between two sets of $d$ nodes. Unlike LoRA, which utilizes an additive, low-rank matrix for weight updates, OFT updates the weights by left-multiplying with an orthogonal matrix. While this block-diagonal arrangement lowers the number of trainable parameters, it is inherently incapable of producing a dense matrix $\bm{R}$. We draw on the structure of butterfly graphs~\cite{cooley1965algorithm}, a network topology famously employed in the Cooley-Tukey algorithm, which is capable of linking $d$ source nodes to $d$ target nodes with only $\mathcal{O}(d\log d)$ connections. In this method, finetuning of a pretrained weight matrix is achieved by representing the new weight matrix as the product of a trainable orthogonal matrix and the fixed, pretrained weight matrix. This explosive growth in scale presents significant obstacles to training foundation models from scratch, rendering such an approach impractical for most researchers. Consequently, we posit that incorporating such structured parameterizations will naturally enforce an inductive bias, which in turn should enhance model generalization and reduce the risk of overfitting. Here, the flow of information initiates with $\bm{B}_1$ and concludes with $\bm{B}_n$. According to the information transmission interpretation, the principal requirements for achieving this goal are: \emph{(i)} \textbf{dense connectivity}, meaning each node at the initial level must have at least one route to every node at the final level, and \emph{(ii)} \textbf{minimum free edges}, implying that the overall count of edges is minimized while maintaining the orthogonality condition. This perspective views BOFT as an optimization over the similarity matrix $\bm{R}$, subject to an orthogonality constraint, thereby closely relating it to the domains of distance metric learning~\cite{xing2002distance} and the theory of bilinear forms~\cite{roman2005advanced}. Butterfly parameterizations also support learning fast linear transforms~\cite{dao2019learning,dao2019kaleidoscope} and have been incorporated for more efficient neural network training~\cite{chen2022pixelated,dao2022monarch}. Our primary contributions are as follows:

\begin{itemize}[leftmargin=*,nosep]
    \item We address parameter efficiency within orthogonal finetuning by developing a new information transmission perspective. \section{Intriguing Insights and Discussions}

\textbf{Expressivity of BOFT}. \item We offer theoretical motivations illustrating how BOFT can achieve notable reductions in the number of tunable parameters without compromising expressivity or the stability of training. The fundamental challenge lies in constructing a dense orthogonal matrix in a parameter-efficient manner. The butterfly network has also been widely adopted in computer architecture~\cite{solihin2015fundamentals,li2011linear} for its effective support of information exchange. Across the board, BOFT demonstrates greater parameter efficiency compared to OFT. We posit that this form of structured inductive bias—arising from butterfly factorization—can be advantageous for model generalization, due to its shared structure with several well-known linear transforms~\cite{dao2019kaleidoscope}, such as the discrete Fourier, sine/cosine, and Hadamard transforms. Because BOFT composes the final orthogonal transformation from several individual orthogonal matrices, its training phase incurs a moderately higher computational overhead relative to OFT. Consequently, it is crucial to develop methods that can adapt such models to new tasks efficiently. When $k\geq 2$ is an integer power of $2$, the \emph{butterfly factor} $\bm{B}^F(k)$ is formally defined as

\begin{equation}\label{eq:but_factor}
\begin{aligned}
    \bm{B}^F(k)=\begin{bmatrix}
    \text{diag}(\bm{d}_1) & \text{diag}(\bm{d}_2) \\
    \text{diag}(\bm{d}_3) & \text{diag}(\bm{d}_4)
    \end{bmatrix}\in\mathbb{R}^{k\times k},
\end{aligned}
\end{equation}

where each $\bm{d}_i \in \mathbb{R}^{\frac{k}{2}},\,\forall i$ represents a vector. The depicted graph emerges from fully unrolling the corresponding matrix multiplications, where each matrix $\bm{B}_i$ characterizes the connections from the $i$-th layer to the $(i+1)$-th layer. \end{document} Overall, the butterfly matrix models a more structured subspace within the orthogonal group (relative to the block-diagonal structure), making BOFT provably more expressive than OFT when block sizes are matched. Within this reframing, synthesizing a dense orthogonal matrix via sparse matrix multiplications resembles relaying information across a grid-like graph topology. \input{rebuttal-results}

\section{Concluding Remarks and Limitations}

In this paper, we introduce Orthogonal Butterfly, a universally applicable, parameter-efficient fine-tuning approach for foundation models, grounded in the butterfly matrix structure. Similar correspondence is observed for decompositions of the form $\bm{R}=\bm{B}_3\bm{B}_2\bm{B}_1$ and $\bm{R}=\bm{B}_2\bm{B}_1$, linking the connectivity patterns in the graphs to the locations of zeros in the factorizations. In contrast, BOFT introduces an alternative fine-tuning method by applying layer-shared transformation matrices to the weights. \textbf{Multiplicative Dropout in BOFT}. & \bm{0} \\
    \bm{0}\!\!\!\!\! The rapid surge in model effectiveness, however, has been accompanied by an exponential increase in parameter count

(\emph{e.g.}, GPT-3 is comprised of approximately 175B parameters). Preserving the spectral norm in this way has been empirically demonstrated to enhance both the stability during training and the model’s capacity to generalize~\cite{miyato2018spectral,yoshida2017spectral}. Ensuring the orthogonality of such a block component $\tilde{\bm{B}}^b(d,2)$ amounts to parameterizing each $2b\times 2b$ block as orthogonal. Numerous prior studies~\cite{dao2019learning,dao2019kaleidoscope,chen2022pixelated,dao2022monarch} have examined sparse training techniques utilizing butterfly parameterizations. Both the butterfly structure and the traditional block-diagonal structure share a common trait: they promote sparsity in the underlying matrices, which serves to lessen the total number of tunable parameters. \section{An Information Transmission View on Orthogonal Finetuning}

We begin by reviewing the foundational aspects of OFT. More precisely, the $(j_1,j_2)$ entry of $\bm{B}_i$ determines the existence of a directed edge from the $j_2$-th node in layer $i$ to the $j_1$-th node in layer $i+1$ (a zero value implies no such edge exists). OFT has shown advantageous properties in terms of generalization and convergence in the context of finetuning text-to-image diffusion models~\cite{qiu2023controlling}. Of these approaches, model finetuning stands out due to its simplicity, efficacy, and, importantly, the fact that it does not introduce extra inference latency. Unlike LoRA’s fixed low-rank design for all added matrices, some recent methods~\cite{zhang2022adaptive,valipour2022dylora,zhang2023increlora} propose dynamically selecting the rank per layer to optimize parameter allocation. Integrating this parameterization into OFT, we present Orthogonal Butterfly (BOFT), a new finetuning strategy that is more parameter-efficient. To ensure that $\bm{B}_5\bm{B}_4\bm{B}_3\bm{B}_2\bm{B}_1$ yields a fully dense matrix, it is necessary that all nodes at the initial layer are capable of passing information to every node in the sixth layer. 
\begin{abstract}
Large foundation models are now pervasive, but training these models from the ground up incurs prohibitive computational costs. In the BOFT approach, we start by initializing every butterfly component to an identity matrix—equivalently, all skew-symmetric matrices for the Cayley transform are initialized as zero matrices. \textbf{Orthogonal finetuning as learning bilinear similarity}. However, due to the multiplicative manner in which BOFT introduces weight changes, standard Dropout is incompatible. Thus, creating a dense orthogonal matrix with standard OFT requires setting $b=d$, whereas BOFT achieves this density for arbitrary values of $b$. To systematically construct such factorizations, we design a graphical framework inspired by information transmission concepts. Experimental results across a variety of domains—including large language models, vision architectures, and text-to-image generators—demonstrate the strong empirical performance of BOFT as a general-purpose fine-tuning strategy. Concretely, if one seeks to express a dense matrix $\bm{R}\in\mathbb{R}^{d\times d}$ as a product of $m$ square matrices, namely, $\thickmuskip=2mu \medmuskip=2mu \bm{R}=\bm{B}_m\bm{B}_{m-1}\cdots\bm{B}_1$, this process may be represented as relaying information across a $d\times (m+1)$ node lattice, as depicted in Figure~\ref{fig:info_trans}. It is customary for parameter-efficient finetuning strategies to begin from the pretrained parameters, thus ensuring the resulting model maintains close fidelity to the original. In contrast, the vanilla OFT with a fully dense orthogonal matrix leverages $\mathcal{O}(d^2)$ parameters, while its block-diagonal variant, using $r$ blocks, involves $\mathcal{O}(bd)$ parameters. \item Drawing inspiration from the butterfly networks in the Cooley-Tukey algorithm, we introduce Orthogonal Butterfly, a novel orthogonal finetuning method characterized by high parameter efficiency, and contextualize it within our information transmission model. Examining a shorter product, say $\bm{R}=\bm{B}_4\bm{B}_3\bm{B}_2\bm{B}_1$, mapping source nodes at layer one to receiving nodes at layer five, reveals, for example, that node $1$ fails to transmit to node $3$, leading to a zero at position $(3,1)$ in the composite's sparsity structure. \end{abstract}

\section{Introduction}

Emerging models such as ChatGPT~\cite{brown2020language,chen2021evaluating} and Stable Diffusion~\cite{rombach2022high} have showcased the extraordinary generalization capabilities inherent in large foundation models. In Figure~\ref{fig:blk_diag}, the block-diagonal configuration of $\bm{R}$ in the standard OFT is illustrated. In particular, the butterfly network underlying Cooley-Tukey introduces a structured approach to sparse matrix factorization termed butterfly factorization. LoRA~\cite{hulora2022}, for example, modifies pretrained weight matrices by incorporating the product of two low-rank matrices, yielding strong performance on natural language tasks. Notably, the butterfly structure is a core component of various efficient linear transforms—such as the discrete Fourier transform and the discrete cosine transform—widely used in practice. Within this framework, designing a dense orthogonal matrix that is parameter-efficient is recast as optimizing an information flow problem on a grid-like graph structure. Starting with the special case where $k=2$, the block-diagonal structure $\tilde{\bm{B}}(d,2)$ can be readily parameterized as an orthogonal matrix, for example by employing the Cayley transform or $2$-dimensional rotation for each block, as outlined in prior works \cite{liu2021orthogonal,qiu2023controlling}. As illustrated in Figure~\ref{fig:info_trans}, a particular matrix factorization can achieve this connectivity with only $10$ edges, which is notably fewer than what a single dense orthogonal matrix would necessitate. Thus, orthogonality may impose the addition or removal of edges for each valid configuration. Likewise, setting $b=d$ yields BOFT($1$,$d$), corresponding to the standard OFT~\cite{qiu2023controlling} featuring an unconstrained orthogonal matrix. Methods for further reducing this training overhead represent a promising direction for future research. The $\bm{R}(m,b)$ structure employed in BOFT commonly selects a structured subset of the orthogonal group, thereby restricting the range of hypotheses and introducing an inductive bias. At first, this appears daunting, since a fully dense orthogonal matrix of size $d$ by $d$ entails $\mathcal{O}(d^2)$ parameters. BOFT tackles the parameter efficiency limitation of OFT by employing butterfly factorization. Thus, balancing expressiveness and structural regularity is crucial for model finetuning generalization. When one chooses $m_1=m_2=\log d$, $b_1=b_2=2$, and $l=d-1$, the matrix $\bm{R}^G(m_1,b_1,m_2,b_2,l)$ can span the entire orthogonal group, a property also referred to as the kaleidoscope hierarchy~\cite{dao2019kaleidoscope}. To resolve this, we propose a specialized multiplicative Dropout for BOFT: Each orthogonal matrix $\bm{R}(m,b)$ comprises $m$ butterfly modules, each of which can be interpreted as a $2b\times 2b$ block-diagonal orthogonal matrix under certain permutations. Similar to LoRA~\cite{hulora2022}, where a Dropout module is included to alleviate overfitting in low-rank adaptation, Dropout~\cite{srivastava2014dropout} can be seamlessly applied in LoRA’s additive update. This construction allows for the butterfly matrix to serve as a parameterization for an orthogonal matrix. The radix-2 Cooley-Tukey algorithm decomposes the $N$-point discrete Fourier transform into a sequence of $\frac{N}{2}$-point transforms, leading to a recursive butterfly pattern which can be described as a product of sparse matrices, commonly termed a butterfly matrix. Following \cite{chen2022pixelated}, we further extend butterfly matrices by introducing the block butterfly component $\tilde{\bm{B}}^b(d,k)$, in which every entry $\bm{d}_i$ is replaced by a $b\times b$ matrix. Notably, after the fine-tuning process, the orthogonal matrices produced by BOFT can be merged directly into the underlying model weights, thereby avoiding extra inference time costs. If one adopts the butterfly configuration, i.e., $m=\log d$ and $b=2$, BOFT’s parameter count becomes $\mathcal{O}(d\log d)$. Unlike earlier work, our focus is on applying butterfly structures to boost the parameter efficiency of OFT. Establishing a fully connected, or dense, linkage between two hierarchical levels requires $\mathcal{O}(d^2)$ connections (that is, forming a dense orthogonal matrix), specifically resulting in $d^2-d$ edges (for $d=4$, this leads to 12 edges). \begin{theorem}[Expressivity of BOFT]\label{thm:exp_boft}
    BOFT possesses greater expressive capacity than OFT for equivalent block sizes. The primary strategies for achieving such efficiency include: \emph{model finetuning}~\cite{hulora2022,zaken2022bitfit,guo2020parameter,raffel2020exploring,qiu2023controlling,zhang2022adaptive,chavan2023one}, where only a subset of existing model parameters are adjusted while maintaining the architecture; \emph{adapter tuning}~\cite{rebuffi2017learning,houlsby2019parameter,pfeiffer2020adapterfusion,lin2020exploring,he2021towards}, which adds additional trainable modules to the original model; and \emph{prompt tuning}~\cite{li2021prefix,lester2021power}, wherein trainable prefix tokens are appended to the model inputs. Although this approach has demonstrated empirical benefits, splitting a neuron's dimensions into $r$ groups according to their indices lacks conceptual justification, underscoring a need for dense orthogonal matrices. The BOFT approach may be reformulated in terms of learning a bilinear similarity, namely $\bm{w}^0_i\bm{R}\bm{x}$, with $\bm{w}^0_i$ representing the $i$-th column (corresponding to a neuron) of the weight matrix $\bm{W}^0$. LoRA, for example, utilizes zero-initialization for its low-rank components. \textbf{Inductive bias and generalization in BOFT}. This transmission perspective unlocks numerous strategies for sparse matrix factorization, each of which curtails parameter count but retains sufficient representational power to produce dense matrices. Since orthogonal matrices demand fewer parameters than unrestricted matrices, this orthogonality constraint induces further dependencies among the edges. These approaches commonly involve directly encoding the weight matrices with butterfly structures and initiating neural network training from the beginning. Choosing $m = \log \frac{2d}{b}$ and $b$ for BOFT($m$,$b$) assigns $\bm{B}^{\frac{b}{2}}(d)$ as $\bm{R}$ and results in a dense orthogonal matrix structure. Additional mathematical aspects of this preservation are discussed in Appendix~\ref{app:sn_property}. &\bm{B}_2(\frac{d}{2})
    \end{bmatrix}= \tilde{\bm{B}}(d,d)\tilde{\bm{B}}(d,\frac{d}{2})\cdots\tilde{\bm{B}}(d,2),
\end{aligned}
\end{equation}

where $\bm{B}_1(\frac{d}{2})$ and $\bm{B}_2(\frac{d}{2})$ represent two butterfly matrices of dimension $\frac{d}{2}$. \item For the first time, we expand the application of orthogonal finetuning to a spectrum of domains beyond its previous use in controllable text-to-image generation~\cite{qiu2023controlling}. This motivation gives rise to the Orthogonal Finetuning~(OFT)~\cite{qiu2023controlling} framework, wherein an identical orthogonal transformation is applied to all neurons within a layer, ensuring that the pairwise angles among neurons remain invariant during finetuning. However, BOFT is not without drawbacks. To investigate this, we conduct a simulation in which a random dense orthogonal matrix~\cite{anderson1987generation} of dimension $512\times 512$ is approximated. Consequently, expressing $\bm{R}$ as a series of factors $\bm{B}_m\bm{B}_{m-1}\cdots\bm{B}_1$ can be interpreted as sequential information propagation governed by the structures of the factor matrices $\bm{B}_i$. In both OFT and BOFT, the finetuning process involves multiplying an orthogonal matrix $\bm{R}$ on the left, resulting in new weights $\bm{R}\bm{U}\bm{\Sigma}\bm{V}^\top$. The requirement of orthogonality places a nuanced restriction on the edges linking consecutive levels. We anticipate that adopting more effective network architectures may further enhance the construction of dense orthogonal matrices. However, the extent to which our orthogonal butterfly matrix can effectively approximate arbitrary orthogonal matrices remains an open question. Although OFT demonstrates strong generalizability, it remains parameter-intensive due to the high-dimensional nature of orthogonal matrices. Each curve sharing the same color represents BOFT with a fixed block size, with the leftmost point corresponding to the error of BOFT($1$,$b$) (i.e., the original block-diagonal OFT with block size $b$). Typical strategies in model finetuning involve using a low learning rate, which serves as a heuristic to keep changes between the pretrained and finetuned models minimal. \section{Orthogonal Parameterization by Butterfly Factorization}

Originating from the Cooley-Tukey algorithm for fast Fourier transform, the butterfly structure enables localized alterations in the frequency domain to induce widespread changes in the spatial domain—an analogy fitting for our scenario, where information at each input node must reach every output node. If orthogonality is maintained in each sparse factor, the entire parameterization becomes highly efficient, utilizing just $\mathcal{O}(d\log d)$ parameters—a substantial improvement over the naive approach. If an entry $R_{ij}$ of matrix $\bm{R}$ is zero, it signifies an absence of transmission from the $j$-th input node to the $i$-th output node; conversely, a nonzero entry permits such transmission. Our empirical results demonstrate that BOFT substantially outperforms existing state-of-the-art techniques, substantiating its advantages in parameter efficiency and generalizability. As seen in Figure~\ref{fig:info_trans}, while the given topology uses $10$ edges to achieve full connectivity, the butterfly configuration can accomplish the same with just $8$ edges. This contrast is illustrated in Figure~\ref{fig:comp_lora}. Visualizing the pattern of nonzero entries in the assembled orthogonal matrix shows why the information transmission framework is particularly effective for analyzing OFT: our primary focus is the nonzero, trainable components of $\bm{R}$, with the actual numeric values being inconsequential. The information transmission perspective provides valuable insights into the sparsity characteristics of the factorization matrices in OFT. Since the sparsity pattern (non-zero entries) of any butterfly component can be rearranged as a permutation of that of $\tilde{\bm{B}}(d,2)$, the orthogonal parameterization can be extended in the same manner to all butterfly components. BOFT enhances the available finetuning parameter space through structural priors, thereby facilitating an improved equilibrium between expressivity and regularity. However, we caution that higher expressivity does not guarantee better fine-tuning outcomes, since full parameterization, despite its expressive universality, can lead to underwhelming empirical results. With the increasing scale and capability of foundation models (\emph{e.g.}, \cite{brown2020language,radford2021learning,rombach2022high,kirillov2023segment}), there has been substantial attention on the problem of adapting these models to new tasks using methods that limit the number of trainable parameters~\cite{treviso2023efficient,ding2023parameter,lialin2023scaling}. This operation leaves the maximal singular value unchanged (i.e., the spectral norm of $\bm{W}^0$ is unaffected). The adoption of a block-diagonal structure enhances parameter-efficiency but imposes a structural limitation: it divides each neuron's dimensions (i.e., the components of each column of $\bm{W}^0$) into $r$ groups, with each group independently transformed by distinct orthogonal matrices. As an illustration, in a $2\times 2$ orthogonal matrix, a zero at position $(1,1)$ necessitates a corresponding zero at $(2,2)$ (that is, the absence of one edge may require another edge to be absent as well). Given that weight matrices possess inherent structure, a more rigorous strategy seeks to preserve the relational geometry—specifically, the pairwise angular relationships between neurons within these matrices. Advancement in this domain, therefore, hinges on the capacity to adapt pre-trained models for novel tasks without resorting to full retraining. Our work departs from convention by assembling a dense orthogonal matrix through the multiplication of several sparse, factorized matrices. Orthogonal finetuning typically produces superior spectral characteristics compared to LoRA, primarily due to its exact preservation of the spectral norm inherent in the original pretrained weight matrix $\bm{W}^0$. Enhancing the parameter efficiency of OFT is thus an important objective.